\documentclass[10pt,a4paper]{amsart}
\usepackage{amsmath,amssymb,amsthm,graphicx}
% For commutative diagrams we use Xy-pic
\usepackage[all]{xy}

% We use 2cell for 2-commutative diagrams.
\xyoption{2cell}
% Theorem environments.
%
\theoremstyle{plain}
\newtheorem{theorem}[subsection]{Theorem}
\newtheorem{proposition}[subsection]{Proposition}
\newtheorem{lemma}[subsection]{Lemma}

\theoremstyle{definition}
\newtheorem{definition}[subsection]{Definition}
\newtheorem{example}[subsection]{Example}
\newtheorem{exercise}[subsection]{Exercise}
\newtheorem{situation}[subsection]{Situation}

\theoremstyle{remark}
\newtheorem{remark}[subsection]{Remark}
\newtheorem{remarks}[subsection]{Remarks}

\numberwithin{equation}{subsection}

% Macros
%
\def\lim{\mathop{\mathrm{lim}}\nolimits}
\def\colim{\mathop{\mathrm{colim}}\nolimits}
\def\Spec{\mathop{\mathrm{Spec}}}
\def\Hom{\mathop{\mathrm{Hom}}\nolimits}
\def\Ext{\mathop{\mathrm{Ext}}\nolimits}
\def\SheafHom{\mathop{\mathcal{H}\!\mathit{om}}\nolimits}
\def\SheafExt{\mathop{\mathcal{E}\!\mathit{xt}}\nolimits}
\def\Sch{\mathit{Sch}}
\def\Mor{\mathop{\mathrm{Mor}}\nolimits}
\def\Ob{\mathop{\mathrm{Ob}}\nolimits}
\def\Sh{\mathop{\mathit{Sh}}\nolimits}
\def\NL{\mathop{N\!L}\nolimits}
\def\CH{\mathop{\mathrm{CH}}\nolimits}
\def\proetale{{pro\text{-}\acute{e}tale}}
\def\etale{{\acute{e}tale}}
\def\QCoh{\mathit{QCoh}}
\def\Ker{\mathop{\mathrm{Ker}}}
\def\Im{\mathop{\mathrm{Im}}}
\def\Coker{\mathop{\mathrm{Coker}}}
\def\Coim{\mathop{\mathrm{Coim}}}

% Boxtimes
%
\DeclareMathSymbol{\boxtimes}{\mathbin}{AMSa}{"02}

%
% Macros for moduli stacks/spaces
%
\def\QCohstack{\mathcal{QC}\!\mathit{oh}}
\def\Cohstack{\mathcal{C}\!\mathit{oh}}
\def\Spacesstack{\mathcal{S}\!\mathit{paces}}
\def\Quotfunctor{\mathrm{Quot}}
\def\Hilbfunctor{\mathrm{Hilb}}
\def\Curvesstack{\mathcal{C}\!\mathit{urves}}
\def\Polarizedstack{\mathcal{P}\!\mathit{olarized}}
\def\Complexesstack{\mathcal{C}\!\mathit{omplexes}}
% \Pic is the operator that assigns to X its picard group, usage \Pic(X)
% \Picardstack_{X/B} denotes the Picard stack of X over B
% \Picardfunctor_{X/B} denotes the Picard functor of X over B
\def\Pic{\mathop{\mathrm{Pic}}\nolimits}
\def\Picardstack{\mathcal{P}\!\mathit{ic}}
\def\Picardfunctor{\mathrm{Pic}}
\def\Deformationcategory{\mathcal{D}\!\mathit{ef}}
\usepackage{cleveref}
\crefname{originallemma}{original lemma}{original lemmas}
\Crefname{originallemma}{Original lemma}{Original lemmas}
\makeatletter
\newlabel{algebra-lemma-hom-from-tensor-product}{{00DE}{source}{}{sourcepointer-1}{}}
\newlabel{algebra-lemma-hom-from-tensor-product@cref}{{[originallemma][0][]00DE}{[1][0][]source}{}{}{}}
\makeatother
\crefname{originallemma}{original lemma}{original lemmas}
\Crefname{originallemma}{Original lemma}{Original lemmas}
\makeatletter
\newlabel{algebra-lemma-evaluation-map-iso-finite-projective}{{0DVB}{source}{}{sourcepointer-2}{}}
\newlabel{algebra-lemma-evaluation-map-iso-finite-projective@cref}{{[originallemma][0][]0DVB}{[1][0][]source}{}{}{}}
\makeatother
\title{Ext of finite projective modules with continuous group action and group cohomology}
\author{The Stacks Project Authors}
\date{stacks-project:etale-cohomology.tex:local-snapshot-20261011; GNU FDL 1.2 or later}
\begin{document}
\maketitle
\noindent\textit{Editorial note.} All delivered mathematical text is the author's own, copied from the pinned Stacks Project chapter file etale-cohomology.tex; no statement or proof step was rewritten, reordered or removed.\par\medskip
\noindent\textit{Editorial note.} The proof is the author's complete argument; the notation referred to in the statement is introduced in the proof as the author indicates.\par\medskip
\noindent\textit{Editorial note.} Every labelled reference occurring inside the counted spans is carried in assembly.external\_references as a non-load-bearing navigation pointer (same-chapter ones at their definition site, cross-chapter ones at the invocation inside the counted span), so no load-bearing mathematics is replaced by a stub. No bibliography citation occurs inside the counted spans.\par\medskip
\noindent\textit{Editorial note.} The checked chapter list for same-name collisions is the set of stacks chapters staged in this repository: chow, cohomology, divisors, etale-cohomology, flat, more-algebra, more-morphisms, morphisms, spaces-more-morphisms, varieties together with this card's chapter etale-cohomology.\par\medskip
\noindent\textit{Editorial note.} The preamble is the chapter's own native preamble excerpt already used for this source: the Xy-pic setup (native preamble lines 14-18) and the native theorem declarations with the mathematical macros (native preamble lines 154-217).\par\medskip
\noindent\textit{Allowed standard prerequisites (editorial).}\begin{itemize}
\item Continuous cohomology of a topological group with coefficients in a module. Same chapter, the cohomology H\textasciicircum{}i(G, M) of a topological group with coefficients in a discrete module, and the standard resolution by continuous functions. Supplies the right hand side of the asserted identification.
\item Ext groups of R-G-modules and finite projectivity of M over R. Same chapter, the Ext groups in the category of R-G-modules and the use of the dual module M\textasciicircum{}vee = Hom\_R(M, R). Supplies the left hand side and the computation by adjunction.
\end{itemize}\medskip
\ifdefined\hypertarget\hypertarget{sourcepointer-1}{}\fi
\noindent\textit{Source reference.} Original lemma 00DE: sources/stacks-etale-cohomology/pinned-original-source.tex, source lines 8923--8923. Navigation pointer invoked at proof line 8923; the counted span contains the pointer only and the cited result is not reproduced inside it, so the pointer is registered as a non-load-bearing navigation stub in the established form for Stacks material. No bibliography citation is involved.\par\medskip
\ifdefined\hypertarget\hypertarget{sourcepointer-2}{}\fi
\noindent\textit{Source reference.} Original lemma 0DVB: sources/stacks-etale-cohomology/pinned-original-source.tex, source lines 8924--8924. Navigation pointer invoked at proof line 8924; the counted span contains the pointer only and the cited result is not reproduced inside it, so the pointer is registered as a non-load-bearing navigation stub in the established form for Stacks material. No bibliography citation is involved.\par\medskip
\begin{lemma}
\label{lemma-ext-modules-hom}
Let $G$ be a topological group. Let $R$ be a ring.
Let $M$, $N$ be $R\text{-}G$-modules. If $M$ is finite projective
as an $R$-module, then
$\text{Ext}^i(M, N) = H^i(G, M^\vee \otimes_R N)$ (for notation
see proof).
\end{lemma}

\begin{proof}
The module $M^\vee = \Hom_R(M, R)$ endowed with the contragredient
action of $G$. Namely $(g \cdot \lambda)(m) = \lambda(g^{-1} \cdot m)$
for $g \in G$, $\lambda \in M^\vee$, $m \in M$. The action of $G$ on
$M^\vee \otimes_R N$ is the diagonal one, i.e., given by
$g \cdot (\lambda \otimes n) = g \cdot \lambda \otimes g \cdot n$.
Note that for a third $R\text{-}G$-module $E$ we have
$\Hom(E, M^\vee \otimes_R N) = \Hom(M \otimes_R E, N)$.
Namely, this is true on the level of $R$-modules by
Algebra, Lemmas \ref{algebra-lemma-hom-from-tensor-product} and
\ref{algebra-lemma-evaluation-map-iso-finite-projective}
and the definitions of $G$-actions are chosen such that it
remains true for $R\text{-}G$-modules. It follows that
$M^\vee \otimes_R N$ is an injective $R\text{-}G$-module
if $N$ is an injective $R\text{-}G$-module. Hence if
$N \to N^\bullet$ is an injective resolution, then
$M^\vee \otimes_R N \to M^\vee \otimes_R N^\bullet$
is an injective resolution. Then
$$
\Hom(M, N^\bullet) = \Hom(R, M^\vee \otimes_R N^\bullet) =
(M^\vee \otimes_R N^\bullet)^G
$$
Since the left hand side computes $\text{Ext}^i(M, N)$ and the right
hand side computes $H^i(G, M^\vee \otimes_R N)$ the proof is complete.
\end{proof}

\end{document}
